\chapter{Planification et architecture}

\section*{Introduction}
\addcontentsline{toc}{section}{Introduction}
Après avoir cadré le besoin métier, ce chapitre traduit la vision produit en exigences, en modèle de données et en architecture technique. Il s'appuie sur les artefacts techniques du projet : documentation d'architecture, schéma de données, couche de routage et services métier \netref{web:project-docs}.

\section{Analyse des besoins et spécification}

\subsection{Identification des acteurs}
Le système distingue quatre acteurs principaux, définis à partir de la spécification fonctionnelle et de la couche de contrôle d'accès. Le tableau \ref{tab:system-actors} synthétise leurs responsabilités.

\begin{table}[H]
  \centering
  \caption{\label{tab:system-actors}Acteurs du système}
  \begin{tabularx}{\textwidth}{>{\bfseries}p{3cm}X}
    \toprule
    Acteur & Responsabilités principales \\
    \midrule
    TA & Gérer le catalogue de CV, créer les offres, affecter les candidats, piloter le matching et les premiers entretiens. \\
    Manager & Évaluer les candidats au stade managérial, conduire les entretiens, produire des rapports et orienter vers les RH. \\
    HR & Finaliser l'évaluation, décider de l'acceptation ou du rejet, générer et envoyer les emails de décision, déclencher l'embauche. \\
    Admin & Superviser les utilisateurs, l'analytique, les logs d'activité, les emails et l'onboarding global. \\
    \bottomrule
  \end{tabularx}
\end{table}

\subsection{Exigences fonctionnelles}
Les exigences fonctionnelles les plus structurantes sont les suivantes :
\begin{itemize}
  \item authentifier les utilisateurs et les rediriger vers un espace adapté à leur rôle ;
  \item importer des CV PDF ou DOCX, extraire leur contenu et enrichir les métadonnées ;
  \item rechercher les profils par mots-clés ou par similarité sémantique ;
  \item créer des postes avec critères \textit{must-have}, \textit{nice-to-have}, séniorité et unité métier ;
  \item affecter les CV à des postes et suivre les candidats dans un pipeline d'états ;
  \item conserver un historique auditable des transitions du pipeline candidat ;
  \item planifier des entretiens, envoyer des invitations et saisir des rapports ;
  \item générer des évaluations IA, des guides d'entretien et des emails ;
  \item exposer un assistant conversationnel capable d'appeler des outils métier ;
  \item soumettre les actions IA modifiant les données à une confirmation explicite ;
  \item fournir des tableaux de bord, des exports et des vues de gouvernance.
\end{itemize}

\subsection{Exigences non fonctionnelles}
Les exigences non fonctionnelles retenues pour la solution sont résumées dans le tableau \ref{tab:nonfunctional-requirements}.

\begin{table}[H]
  \centering
  \caption{\label{tab:nonfunctional-requirements}Exigences non fonctionnelles clés}
  \begin{tabularx}{\textwidth}{>{\bfseries}p{3.2cm}X}
    \toprule
    Axe & Traduction dans la solution \\
    \midrule
    Sécurité & Authentification \sourcecode{better-auth}, contrôle d'accès par rôle sur les routes, les actions serveur et les outils IA \netref{web:better-auth-docs}. \\
    Validation & Validation \sourcecode{Zod} sur les entrées utilisateur et les arguments d'outils \netref{web:zod-docs}. \\
    Performance & Sélection explicite des colonnes, limites de requêtes, embeddings asynchrones, streaming SSE côté chat. \\
    Évolutivité & Architecture \textit{feature-driven}, services métier séparés, modules IA découplés, stockage \sourcecode{pgvector} \netref{web:pgvector-docs}. \\
    Traçabilité & Logs d'activité, email logs, notifications, conversations, messages, historique d'étapes candidat et confirmations d'actions IA stockés en base. \\
    Fiabilité IA & Guardrails sur le nombre d'étapes agentiques, les tokens, le timeout, les échecs d'outils et la confirmation des actions modifiant les données. \\
    \bottomrule
  \end{tabularx}
\end{table}

\section{Étude analytique}

\subsection{Diagramme global des cas d'utilisation}
La figure \ref{fig:global-use-case} répartit les capacités entre les quatre acteurs métier réels : Talent Acquisition, Manager métier, Ressources humaines et Administrateur. L'authentification constitue un prérequis commun ; elle n'est pas répétée sous la forme d'un acteur supplémentaire, afin de préserver une lecture directe des responsabilités.

\begin{figure}[H]
  \centering
  \diagraminput{figures/diagrams/global-use-case}
  \caption{\label{fig:global-use-case}Cas d'utilisation de la plateforme \projectname{}, structurés par rôle}
\end{figure}

\subsection{Processus métier cible}
Le diagramme de cas d'utilisation identifie les capacités du système, mais il ne suffit pas à montrer la circulation réelle d'une candidature entre les rôles. Pour répondre à cette limite, la figure \ref{fig:business-bpmn} représente le processus cible dans un pool découpé en couloirs horizontaux : Talent Acquisition initialise la candidature, le système IA assiste la présélection, le Manager évalue la compétence opérationnelle, puis les Ressources humaines arbitrent la décision finale. Chaque passerelle exclusive conduit soit à l'étape suivante, soit à une notification de rejet suivie de la clôture de la candidature.

\begin{figure}[H]
  \centering
  \diagraminput{figures/diagrams/recruitment-bpmn}
  \caption{\label{fig:business-bpmn}Processus BPMN cible du recrutement avec pool et couloirs de responsabilité}
\end{figure}

Cette lecture métier justifie ensuite les choix de conception : chaque étape doit posséder un responsable, un état observable, une trace et une règle de transition. Le pipeline candidat et l'historique des changements ne sont donc pas des détails techniques ; ils matérialisent le besoin de gouvernance visible dans le processus.

\subsection{Diagramme de classes global}
Le schéma Drizzle du projet est reformulé en diagramme de classes UML afin de montrer les classes métier, leurs attributs principaux et leurs relations structurantes. La figure \ref{fig:domain-model} distingue explicitement le CV de la candidature créée lors de son affectation à une offre. Elle présente aussi l'historique des étapes, les actions agentiques en attente et la dépendance directe du rapport à l'entretien, conformément aux clés étrangères du schéma.

\begin{figure}[H]
  \centering
  \diagraminput{figures/diagrams/global-domain-model}
  \caption{\label{fig:domain-model}Modèle de domaine et relations structurantes du recrutement}
\end{figure}

\subsection{Vue d'ensemble du Product Backlog}
Le \textit{Product Backlog} constitue la passerelle entre le besoin métier décrit au chapitre précédent et la réalisation par sprints présentée dans la suite du rapport. Il ne s'agit pas seulement d'une liste de fonctionnalités : il ordonne les capacités du produit selon leur valeur pour les équipes de recrutement, leurs dépendances techniques et leur rôle dans la constitution progressive des données exploitables par l'IA, comme le montre le tableau \ref{tab:product-backlog}.

\paragraph{Logique de priorisation.}
La priorité a été donnée aux éléments qui débloquent les autres modules. Le socle d'accès arrive en premier parce qu'aucune donnée RH ne peut être exposée sans authentification et contrôle de rôle. Le catalogue de CV et les offres viennent ensuite, car ils produisent les objets métier fondamentaux. Le pipeline candidat transforme ces objets en processus opérationnel. L'IA intervient lorsque suffisamment de données structurées existent pour produire des scores, des synthèses et des réponses contextualisées. Enfin, la gouvernance consolide les traces, les notifications et les indicateurs nécessaires au pilotage.

\begin{table}[H]
  \centering
  \caption{\label{tab:product-backlog}Product Backlog synthétique du projet}
  \begin{tabularx}{\textwidth}{>{\bfseries}p{1.2cm}p{3cm}p{3.2cm}X}
    \toprule
    ID & Epic produit & Valeur métier & Items principaux \\
    \midrule
    PB1 & Accès et identité & Sécuriser l'entrée dans la plateforme et orienter chaque utilisateur vers son espace. & Connexion, redirection par rôle, protection des espaces, shell de dashboard. \\
    PB2 & Catalogue de CV & Centraliser les profils et transformer des documents hétérogènes en données consultables. & Import PDF/DOCX, extraction de texte, métadonnées, recherche, export, détection de doublons. \\
    PB3 & Offres d'emploi & Structurer les besoins de recrutement pour préparer le matching et les décisions. & Création de poste, critères \textit{must-have}, critères \textit{nice-to-have}, séniorité, fermeture contrôlée. \\
    PB4 & Pipeline candidat & Faire progresser un profil affecté dans un workflow multi-rôles traçable. & Affectation, états candidat, responsabilités TA/Manager/HR, planification, rapports d'entretien, historique immuable des transitions. \\
    PB5 & Intelligence IA & Assister les recruteurs par des analyses contextualisées et contrôlées. & Screening, matching hybride, guides d'entretien, autopilot, RAG, chat analytique outillé, confirmation des actions modifiant les données. \\
    PB6 & Gouvernance & Donner de la visibilité sur l'activité, les communications et l'onboarding. & Notifications, logs d'activité, logs emails, analytics, exports, centre de gouvernance, suivi d'intégration. \\
    \bottomrule
  \end{tabularx}
\end{table}

Ce backlog permet une lecture produit du projet : chaque epic correspond à une capacité visible par un utilisateur final, mais aussi à une brique de données réutilisable par les modules suivants. Par exemple, les offres structurées alimentent le matching ; les rapports d'entretien enrichissent la traçabilité ; les logs et notifications rendent possible une supervision administrateur crédible.

\subsection{Découpage en sprints et objectifs}
Le backlog est découpé en cinq sprints de réalisation. Chaque sprint possède un objectif explicite, afin que la progression du rapport reste lisible et que chaque chapitre d'implémentation puisse être évalué par rapport à une intention produit précise. Le tableau \ref{tab:sprint-objectives} formalise ce découpage.

\begin{table}[H]
  \centering
  \caption{\label{tab:sprint-objectives}Objectifs des sprints de réalisation}
  \begin{tabularx}{\textwidth}{>{\bfseries}p{1.7cm}p{3.2cm}X}
    \toprule
    Sprint & Module & Objectif principal \\
    \midrule
    Sprint 1 & Accès et identité & Mettre en place le point d'entrée sécurisé du produit : authentification, redirection par rôle et cadre de navigation commun. \\
    Sprint 2 & CV et offres & Construire la base documentaire et métier du recrutement : importer les CV, structurer les offres et préparer les traitements de recherche et de matching. \\
    Sprint 3 & Pipeline et entretiens & Transformer une affectation CV-offre en suivi candidat complet, avec transitions d'état, historique auditable, responsabilités, planification et rapports. \\
    Sprint 4 & Intelligence IA & Ajouter une assistance IA réellement branchée sur les données internes : screening, recherche hybride, guides, autopilot, chat analytique et confirmation des actions modifiant les données. \\
    Sprint 5 & Gouvernance & Superviser le produit en exploitation : notifications, emails, activité, analytics, exports, centre de gouvernance et onboarding. \\
    \bottomrule
  \end{tabularx}
\end{table}

\begin{figure}[H]
  \centering
  \diagraminput{figures/diagrams/sprint-roadmap}
  \caption{\label{fig:sprint-roadmap}Trajectoire incrémentale des sprints de réalisation}
\end{figure}

Cette planification montre une trajectoire progressive, résumée par la figure \ref{fig:sprint-roadmap} : d'abord sécuriser l'accès, ensuite créer les données de référence, puis orchestrer le workflow métier, enrichir les décisions par l'IA et terminer par la supervision. Elle évite de présenter les sprints comme des blocs indépendants ; chacun prépare explicitement le suivant.

\section{Architecture du projet}

\subsection{Architecture externe}
D'un point de vue global, la plateforme peut être comprise comme un système organisé en couches. Cette lecture évite de réduire l'architecture à une liste de fichiers ou de technologies : elle met plutôt en évidence les responsabilités principales du produit et la manière dont les données circulent entre les utilisateurs, la logique métier, les données persistantes et les services externes.
\begin{itemize}
  \item la couche présentation regroupe les espaces utilisés par les rôles \textbf{TA}, \textbf{Manager}, \textbf{RH} et \textbf{Admin} ;
  \item la couche d'accès et d'orchestration contrôle l'identité, les rôles, les redirections et la coordination des actions ;
  \item la couche métier porte les règles de recrutement : catalogue de CV, offres, candidats, entretiens, décisions, IA et gouvernance ;
  \item la couche données et services partagés conserve les informations, les historiques, les notifications, les exports et les appels vers les services spécialisés.
\end{itemize}

Le flux typique est le suivant : un utilisateur agit depuis son espace métier, la demande est contrôlée selon son rôle, la règle métier correspondante est appliquée, puis les données sont consultées ou mises à jour. Lorsqu'un cas d'usage l'exige, la plateforme sollicite un service spécialisé comme l'assistance IA ou l'envoi d'email, tout en conservant la décision et les preuves dans le périmètre applicatif.

\begin{figure}[H]
  \centering
  \diagraminput{figures/diagrams/layered-architecture}
  \caption{\label{fig:layered-architecture}Vue globale de l'architecture par couches}
\end{figure}

\subsection{Architecture applicative globale}
L'architecture adoptée n'est pas une architecture microservices. Le projet repose sur un monolithe moderne organisé par fonctionnalités. Ce choix signifie que la plateforme reste livrée comme une application cohérente, tout en séparant clairement les capacités métier : accès, catalogue de CV, offres, pipeline candidat, entretiens, intelligence assistive, gouvernance et administration.

\begin{figure}[H]
  \centering
  \diagraminput{figures/diagrams/monolithic-architecture}
  \caption{\label{fig:monolithic-architecture}Architecture monolithique modulaire de la plateforme}
\end{figure}

Ce choix est adapté au contexte du projet, car les modules partagent les mêmes utilisateurs, les mêmes rôles, les mêmes candidats, les mêmes règles d'accès et la même base de preuves. Les séparer artificiellement en microservices aurait augmenté la complexité de synchronisation, de sécurité et de déploiement sans créer de bénéfice métier immédiat.

\subsection{Architecture interne}
Les frontières internes se lisent donc en termes de responsabilités :
\begin{itemize}
  \item la présentation porte les écrans, formulaires, tableaux de bord et interactions utilisateur ;
  \item l'accès contrôle l'authentification, les rôles et l'orientation vers les bons espaces ;
  \item les modules métier concentrent les règles propres au recrutement ;
  \item les données persistantes conservent les profils, offres, candidats, entretiens, historiques et preuves ;
  \item les services transverses apportent les capacités partagées : notifications, exports, assistance IA, recherche et audit.
\end{itemize}

\paragraph{Tableau de correspondance besoin / module / solution technique.}
La cohérence de l'architecture se comprend mieux lorsqu'on relie directement les besoins métier aux modules fonctionnels et aux réponses qui y répondent, comme le montre le tableau \ref{tab:business-architecture-correspondence}.

\begin{table}[H]
  \centering
  \caption{\label{tab:business-architecture-correspondence}Correspondance entre besoins métier et réponses architecturales}
  \begin{tabularx}{\textwidth}{>{\bfseries}p{3.1cm}p{3.1cm}X}
    \toprule
    Besoin métier & Module principal & Réponse retenue \\
    \midrule
    Sécuriser l'accès et séparer les responsabilités & Accès et identité & Authentification par rôles, espaces protégés et cadre de navigation mutualisé. \\
    Transformer des CV hétérogènes en profils exploitables & Catalogue de CV & Pipeline d'ingestion, extraction structurée, enrichissement progressif et recherche. \\
    Relier les profils aux besoins de recrutement & Offres, candidats et pipeline & Offres structurées, transitions d'état explicites, historique d'étapes et assignations multi-rôles. \\
    Retrouver des profils au-delà des mots-clés & Intelligence IA & Recherche assistée combinant signaux métier, similarité sémantique et justification des résultats. \\
    Interroger le système en langage naturel & Chat analytique & Assistant conversationnel outillé, réponses sourcées, confirmations et traces d'exécution. \\
    Auditer et piloter l'activité & Gouvernance & Notifications, logs, centre de gouvernance, analytics, exports et suivi détaillé de l'onboarding. \\
    \bottomrule
  \end{tabularx}
\end{table}

Ce tableau montre que l'architecture n'est pas le résultat d'un empilement de technologies. Chaque décision répond à un besoin métier identifiable et s'insère dans un module fonctionnel précis.

\paragraph{Justification du choix feature-driven.}
Le choix \textit{feature-driven} est pertinent parce que le produit s'organise naturellement autour de tranches métier cohérentes : accès, catalogue de CV, offres, pipeline candidat, IA et gouvernance. Dans ce contexte, regrouper les éléments d'un même module autour de sa capacité métier réduit la dispersion et rend plus lisible la responsabilité de chaque partie.

\paragraph{Justification du monolithe moderne.}
Le projet n'appelle pas naturellement une architecture microservices. Les modules partagent fortement les mêmes entités, la même authentification, le même stockage et les mêmes règles de gouvernance. Un monolithe moderne permet donc de garder une base de vérité unifiée pour les rôles, les candidats, les entretiens, le chat et les traces d'activité. Il réduit également la complexité de déploiement, de cohérence transactionnelle et de sécurité, tout en restant compatible avec des capacités avancées comme les réponses progressives, l'orchestration de l'assistance IA et la gouvernance des actions sensibles.

\subsection{Architecture du sous-système IA}
Le sous-système IA articule plusieurs briques : génération de descriptions de poste, extraction de CV, matching, génération de guides d'entretien, débriefing, assistant conversationnel et pipeline RAG. L'agent statistique est exposé via l'API du chat analytique, qui assemble le contexte, filtre les outils selon le rôle, diffuse la réponse en SSE et restitue les preuves sous forme de sources, graphiques, cartes typées et trace d'exécution par phases.
La figure \ref{fig:ai-orchestration} illustre cette orchestration entre le chat, les outils métier, les garde-fous et les preuves restituées.

\begin{figure}[H]
  \centering
  \diagraminput{figures/diagrams/ai-orchestration}
  \caption{\label{fig:ai-orchestration}Architecture du sous-système agentique IA et RAG}
\end{figure}

\paragraph{Justification de \sourcecode{pgvector}.}
L'usage de \sourcecode{pgvector} permet de conserver les embeddings dans la même base que les données métier. Ce choix s'appuie sur la documentation officielle de l'extension \sourcecode{pgvector} \netref{web:pgvector-docs} et évite de déporter la recherche vectorielle vers une infrastructure séparée. Il maintient une forte proximité entre profils, chunks, candidats et services de retrieval. Dans un produit de recrutement où les données doivent rester gouvernées, traçables et liées aux mêmes règles d'accès, cette unification simplifie l'architecture sans empêcher les usages avancés de similarité sémantique.

\paragraph{Justification du retrieval hybride.}
Le retrieval hybride est retenu parce qu'aucun signal unique n'est suffisant dans un contexte RH. Le principe RAG renvoie aux travaux de Lewis et al. sur la génération augmentée par récupération \netref{web:rag-lewis}, tandis que la recherche vectorielle s'inscrit dans les approches de similarité à grande échelle mentionnées en bibliographie \netref{web:vector-johnson}. La recherche lexicale reste utile pour les intitulés exacts, les technologies nommées ou les mots-clés métier, tandis que la recherche vectorielle capte des proximités sémantiques moins littérales. Leur combinaison via RRF, complétée par un reranking, permet de limiter à la fois les faux négatifs liés au vocabulaire et les rapprochements trop flous liés à la seule similarité vectorielle.

\paragraph{Justification de l'agent outillé.}
Le choix d'un agent outillé répond à une limite claire des chatbots génériques : produire une réponse fluide ne suffit pas lorsqu'il faut interroger des données sensibles, filtrer les accès et agir sur des objets métier réels. En s'appuyant sur un registre d'outils explicitement validés, l'agent peut raisonner sur plusieurs étapes tout en restant borné par les opérations autorisées. Il devient ainsi une interface d'accès au système, et non un simple générateur de texte détaché du contexte applicatif.

\paragraph{Sécurité et fiabilité IA.}
Cette couche est encadrée par plusieurs mécanismes de sécurité et de fiabilité :
\begin{itemize}
  \item RBAC, qui restreint l'accès aux outils, aux vues et aux données selon le rôle utilisateur ;
  \item validation Zod, qui borne les entrées des outils et évite qu'un appel mal formé ne soit exécuté implicitement ;
  \item confirmation explicite, qui transforme les outils agentiques modifiant les données en actions en attente avant exécution ;
  \item \textit{grounding}, qui interdit à la réponse finale de citer des candidats non présents dans les résultats réellement accessibles ;
  \item cartes typées et trace par phases, qui transforment les sorties d'outils en métriques lisibles, liens profonds, actions de suivi et timeline d'exécution auditable ;
  \item isolation de cache et de périmètre, qui évite qu'une recherche ou un contexte d'un utilisateur ne contamine celui d'un autre ;
  \item garde-fous d'exécution, comme les limites d'étapes agentiques, de tokens, de timeout et d'échecs consécutifs ;
  \item suite d'évaluation agentique, qui vérifie de manière déterministe les politiques de grounding, de confirmation, de cartes de réponse, de traces d'outils et de masquage des données sensibles.
\end{itemize}

Ces garde-fous sont essentiels, car la crédibilité du projet repose précisément sur une IA utile, gouvernée et mesurable, et non sur une automatisation opaque.

\section{Environnement de développement et technique}

\subsection{Outils de gestion et de documentation}
L'environnement de travail est présenté à travers les artefacts utiles à la compréhension du projet : documents de conception, schéma de données, modules fonctionnels, diagrammes, captures, scénarios de validation et rapports d'évaluation.

Cette présentation garde l'analyse au niveau académique attendu : ce qui importe n'est pas l'organisation locale des fichiers, mais la manière dont les artefacts soutiennent les choix de conception et les preuves de réalisation.

\subsection{Environnement d'exécution et de développement}
Le projet s'appuie sur un environnement web moderne permettant de développer une application full-stack, de gérer la persistance, de valider les entrées, d'exécuter les tests et de produire les évaluations nécessaires. La section précise le rôle de l'outillage dans la réalisation : construire l'interface, organiser la logique métier, stocker les données, automatiser certaines validations et mesurer les modules d'intelligence.

Ce choix d'outillage simplifie le cycle de réalisation : une même base technique permet de construire les écrans, d'orchestrer les traitements métier, de contrôler les données et de vérifier les comportements critiques sans disperser le projet entre plusieurs environnements incohérents.

\subsection{Outils front-end}
Le front-end du projet repose sur les technologies suivantes :
\begin{itemize}
  \item \sourcecode{Next.js 16} avec App Router pour le rendu, les layouts et l'orchestration des routes \netref{web:nextjs-docs} ;
  \item \sourcecode{React} pour la composition de l'interface \netref{web:react-docs} ;
  \item \sourcecode{TypeScript} en mode strict pour la sécurité de typage \netref{web:typescript-docs} ;
  \item \sourcecode{Tailwind CSS v4} et \sourcecode{shadcn/ui} pour les composants et le design system \netref{web:tailwind-docs}\netref{web:shadcn-docs} ;
  \item \sourcecode{@tabler/icons-react} pour l'iconographie ;
  \item \sourcecode{React Hook Form} pour les formulaires complexes ;
  \item \sourcecode{TanStack Query} lorsque le cache client est utile ;
  \item \sourcecode{Recharts} pour les visualisations analytiques ;
  \item \sourcecode{Framer Motion} pour certaines micro-interactions.
\end{itemize}

Ce socle front-end répond au besoin de construire des interfaces riches mais maintenables. Les dashboards TA, Manager, HR et Admin doivent rester lisibles tout en exposant des fonctionnalités avancées comme l'analytique, les assistants IA et les vues de pipeline.

\subsection{Outils back-end}
Le back-end est intégré dans la même application full-stack, mais organisé de manière stricte :
\begin{itemize}
  \item les routes \sourcecode{Next.js} gèrent les points d'entrée HTTP et le streaming \sourcecode{SSE} ;
  \item les Server Actions servent de couche de mutation contrôlée ;
  \item les services métier concentrent la logique de traitement et les accès aux données ;
  \item \sourcecode{better-auth} gère l'authentification, les sessions et les rôles \netref{web:better-auth-docs} ;
  \item \sourcecode{Zod} valide les entrées utilisateur et les arguments des outils agentiques \netref{web:zod-docs} ;
  \item \sourcecode{Nodemailer} assure l'envoi des emails d'entretien et de décision ;
  \item le SDK compatible \sourcecode{OpenAI} est utilisé avec des fournisseurs IA externes pour les appels LLM et embeddings ;
  \item \sourcecode{Vitest} couvre les tests automatisés observables du projet \netref{web:vitest-docs}.
\end{itemize}

Cette organisation permet de mélanger logique métier classique et logique IA sans éclater le projet en sous-systèmes difficiles à maintenir.

\subsection{Architecture de la base de données}
La persistance s'appuie sur \sourcecode{Neon PostgreSQL} et \sourcecode{Drizzle ORM} \netref{web:neon-docs}\netref{web:postgres-docs}\netref{web:drizzle-docs}. Ce choix permet de garder dans la même base :
\begin{itemize}
  \item les données métier du recrutement : utilisateurs, offres, CV, candidats, entretiens et rapports ;
  \item les données conversationnelles du chat IA ;
  \item les journaux d'activité, notifications et logs d'emails ;
  \item l'historique des transitions candidat via \sourcecode{candidate\_stage\_history} ;
  \item les confirmations d'actions IA via \sourcecode{pending\_agent\_actions} ;
  \item les tâches d'onboarding ;
  \item les embeddings vectoriels via \sourcecode{pgvector}.
\end{itemize}

L'usage de \sourcecode{pgvector} évite l'introduction d'une base vectorielle externe \netref{web:pgvector-docs}. Les embeddings globaux des CV sont stockés dans le catalogue, tandis que les embeddings orientés RAG sont stockés dans les chunks. Ces chunks sont aussi indexés en recherche plein texte, ce qui rend possible une recherche hybride combinant récupération vectorielle et récupération lexicale.
Les tables \sourcecode{candidate\_stage\_history} et \sourcecode{pending\_agent\_actions} renforcent la gouvernance du modèle. La première conserve les mouvements du pipeline avec l'étape précédente, la nouvelle étape, l'acteur, la raison, la source et l'horodatage. La seconde matérialise les actions IA sensibles avant exécution, avec leur statut, leur expiration et leur résultat d'exécution. Ces deux tables transforment les opérations critiques en preuves consultables et exportables.

\subsection{Justification synthétique des choix technologiques}
Les technologies retenues ne sont pas uniquement des préférences de développement. Elles répondent à des contraintes mesurables du projet : rapidité d'itération, typage strict, gouvernance des données RH, capacité de streaming pour l'agent, recherche hybride et maintien d'un périmètre déployable par une petite équipe. Le tableau \ref{tab:tech-choice-justification} synthétise cette logique de décision.

\begin{table}[H]
  \centering
  \scriptsize
  \caption{\label{tab:tech-choice-justification}Justification des principaux choix technologiques}
  \begin{adjustbox}{max width=\linewidth,center}
  \begin{tabularx}{\textwidth}{>{\bfseries\raggedright\arraybackslash}p{2.5cm}>{\raggedright\arraybackslash}p{2.8cm}>{\raggedright\arraybackslash}X>{\raggedright\arraybackslash}X}
    \toprule
    Choix retenu & Alternatives possibles & Justification pour le projet & Risque maîtrisé \\
    \midrule
    Next.js 16 App Router & SPA séparée, API backend isolée & Unifie rendu, routes, actions serveur et streaming SSE dans un même produit full-stack. & Réduit la dispersion entre front-end, API et logique d'orchestration. \\
    TypeScript strict & JavaScript simple & Rend explicites les contrats entre composants, services, actions et outils agentiques. & Diminue les erreurs d'intégration dans les workflows multi-rôles. \\
    Bun & npm, yarn, pnpm, Node scripts dispersés & Accélère les scripts locaux et garde une convention unique pour développement, tests, seed et évaluation. & Évite la fragmentation de l'outillage. \\
    Architecture feature-driven & Découpage par couches techniques ou microservices & Aligne chaque module sur une capacité métier : accès, catalogue de CV, jobs, pipeline, IA, gouvernance. & Évite un microservice prématuré et maintient une cohérence transactionnelle. \\
    Better Auth + RBAC & Authentification maison & Fournit sessions et rôles sans réimplémenter une couche sensible. & Limite les failles liées à une authentification artisanale. \\
    Zod & Validation implicite ou types compile-time seuls & Valide les entrées utilisateur et les arguments d'outils à l'exécution. & Empêche l'agent ou les formulaires de déclencher des mutations mal formées. \\
    Neon PostgreSQL + Drizzle & Base NoSQL ou ORM plus opaque & Garde données RH, logs, conversations, états et embeddings dans un modèle relationnel typé. & Préserve traçabilité, cohérence et requêtage analytique. \\
    pgvector & Base vectorielle externe & Ajoute la similarité sémantique sans séparer les embeddings des règles métier et d'accès. & Réduit la complexité d'infrastructure et les risques de divergence des droits. \\
    Tailwind CSS + shadcn/ui & CSS ad hoc ou bibliothèque fermée & Combine rapidité d'interface, accessibilité des primitives et personnalisation visuelle. & Évite de reconstruire des composants interactifs sensibles comme les dialogues. \\
    \bottomrule
  \end{tabularx}
  \end{adjustbox}
\end{table}

Cette matrice rend explicite le critère principal de conception : préférer une architecture intégrée, typée et gouvernée à une architecture plus distribuée mais plus coûteuse à fiabiliser. Ce choix est cohérent avec la nature du projet : une plateforme de recrutement interne où la cohérence des données, la sécurité des rôles et la traçabilité des décisions priment sur une scalabilité horizontale prématurée.

\section*{Conclusion}
\addcontentsline{toc}{section}{Conclusion}
Le produit est désormais défini à trois niveaux : les acteurs, les besoins et la structure technique qui permet d'y répondre. Les chapitres suivants décrivent la réalisation concrète de cette architecture, en commençant par le socle d'accès, d'authentification et d'orientation par rôle.
